\section{FluxOS: orchestration, placement, and the application lifecycle}
\label{sec:fluxos}

\FluxOS is the Node.js daemon that every \FluxNode operator runs. It exposes the node's HTTP/HTTPS
API, keeps a gossip mesh open to other \FluxOS instances, and decides — locally, with no message
sent to any other component asking permission — whether to pull an image, start a container, leave
it running, redeploy it, or tear it down. There is no scheduler in the sense a reader from
conventional cluster orchestration would expect: no process anywhere holds a global view of node
load and assigns work to relieve it. Instead, every node runs the same loop over the same
gossiped state and reaches its own conclusion about which of the network's under-replicated
applications to attempt next. What keeps two thousand independent conclusions from producing three
thousand copies of the same container is not agreement — it is a fixed wait, a broadcast, and a
rule for breaking ties by whoever broadcast first. That mechanism, its guarantees, and its failure
mode are the center of this section. Everything below is \shipped and cited to the deployed
\FluxOS source unless a claim is explicitly marked otherwise; where the evidence gathered for this
paper itself records a code path as not fully traced, this section says so rather than smoothing
the gap into confident prose. Line numbers in \texttt{flux/\ldots} citations are read against the
organisation repository's working checkout at commit \texttt{933614b1c}, one unpublished commit
ahead of \texttt{master} --- the flagged \specv{9} schema gate described in \S\ref{sec:appspec-v9},
which also appends a price-table row at height \num{2950000} with unchanged rates; where a cited
construct exists only on that branch, the citation says so.

\subsection{Process architecture}
\label{sec:fluxos-process}

\subsubsection{Processes and ports}
\FluxOS starts as two entry points: \texttt{apiServer.js} (the HTTP/HTTPS API) and
\texttt{homeServer.js} (a static UI server) \srccode{flux/\allowbreak{}apiServer.js}{285}. \texttt{apiServer.js}
calls \texttt{serviceManager.\allowbreak{}startFluxFunctions()} only after both listeners are up
\srccode{flux/\allowbreak{}apiServer.js}{285}, so the boot sequence described below begins after the ports are
already bound.

\begin{table}[H]
\centering
\caption{\FluxOS process ports and adjacent service ports.}
\begin{tabular}{l l l}
\toprule
Listener & Value & Source \\
\midrule
API port (default) & 16127 & \srccode{flux/\allowbreak{}ZelBack/\allowbreak{}config/\allowbreak{}default.js}{28} \\
Allowed API ports (boot aborts otherwise) & \{16127,16137,16147,16157,16167,16177,16187,16197\} &
  \srccode{flux/\allowbreak{}ZelBack/\allowbreak{}config/\allowbreak{}default.js}{26} \\
API HTTPS port & apiPort $+$ 1 & \srccode{flux/\allowbreak{}apiServer.js}{38} \\
Home UI port & apiPort $-$ 1 & \srccode{flux/\allowbreak{}homeServer.js}{19} \\
Syncthing (by convention) & apiPort $+$ 2, absolute default 8384 & \srcdoc{flux/ZelBack/config/default.js:405} \\
MongoDB & 127.0.0.1:27017 & \srccode{flux/\allowbreak{}ZelBack/\allowbreak{}config/\allowbreak{}default.js}{31--32} \\
\fluxbench RPC / P2P port & 16224 / 16225 (testnet 26224 / 26225) & \srccode{flux/\allowbreak{}ZelBack/\allowbreak{}config/\allowbreak{}default.js}{102--106} \\
\fluxd P2P / RPC / ZMQ port & 16125 / 16124 / 16123 (testnet 26125 / 26124) & \srccode{flux/\allowbreak{}ZelBack/\allowbreak{}config/\allowbreak{}default.js}{109--115} \\
\bottomrule
\end{tabular}
\end{table}

The HTTPS listener uses a self-signed certificate generated on first boot by
\texttt{helpers/\allowbreak{}createSSLcert.\allowbreak{}sh} if \texttt{certs/v1.key} is absent
\srccode{flux/\allowbreak{}apiServer.js}{248--258}. Boot hard-blocks on \texttt{dbHelper.\allowbreak{}waitForMongo()} and
\texttt{dockerService.\allowbreak{}waitForDocker()} before any other service starts; the code's own comment
reads "hard dependencies — nothing starts until these are confirmed"
\srcdoc{flux/ZelBack/src/services/serviceManager.js:132} \srccode{flux/\allowbreak{}ZelBack/\allowbreak{}src/\allowbreak{}services/\allowbreak{}serviceManager.js}{133--134}.

\subsubsection{Boot sequence}
\texttt{serviceManager.\allowbreak{}startFluxFunctions} runs a long, mostly-\texttt{await}ed sequence
\srccode{flux/\allowbreak{}ZelBack/\allowbreak{}src/\allowbreak{}services/\allowbreak{}serviceManager.js}{126--572}. In order: (1) reject an
unsupported API port and exit \srccode{flux/\allowbreak{}ZelBack/\allowbreak{}src/\allowbreak{}services/\allowbreak{}serviceManager.js}{126--129}; (2)
start the enterprise node/owner map sync from \texttt{helpers/\allowbreak{}enterprisenodes.\allowbreak{}json}, then resync
from GitHub every 6\,h, with fetch failure swallowed by a \texttt{.catch} handler so boot proceeds
on stale data \srccode{flux/\allowbreak{}ZelBack/\allowbreak{}src/\allowbreak{}services/\allowbreak{}serviceManager.js}{130--136}; (3) hard-block on
Mongo and Docker availability (above); (4) start UPnP firewall maintenance if a router IP is
configured \srccode{flux/\allowbreak{}ZelBack/\allowbreak{}src/\allowbreak{}services/\allowbreak{}serviceManager.js}{138--153}; (5) start
\texttt{flux\allowbreak{}Node\allowbreak{}Service}, repair Docker log corruption, and start \texttt{systemService.\allowbreak{}monitorSystem()}
\srccode{flux/\allowbreak{}ZelBack/\allowbreak{}src/\allowbreak{}services/\allowbreak{}serviceManager.js}{158--163}; (6) drop the
\texttt{loggedusers}\slash\texttt{activeloginphrases}\slash\texttt{activesignatures} collections — every
restart invalidates all sessions — and (re)create TTL indexes: \texttt{loggedusers} 14\,d,
\texttt{activeloginphrases}\slash\texttt{activesignatures} 900\,s, \texttt{activepaymentrequests}
3{,}600\,s, \texttt{completedpayments} 7\,d, \texttt{apptamperingevents} 30\,d
\srccode{flux/\allowbreak{}ZelBack/\allowbreak{}src/\allowbreak{}services/\allowbreak{}serviceManager.js}{165--215}; (7) run tamper-reboot detection
and merge pre-unique-index duplicate runtime-state documents
\srccode{flux/\allowbreak{}ZelBack/\allowbreak{}src/\allowbreak{}services/\allowbreak{}serviceManager.js}{216--220}; (8) prepare global-apps indexes,
including TTL-by-\texttt{expireAt} on the location, event, and installing-broadcast collections
\srccode{flux/\allowbreak{}ZelBack/\allowbreak{}src/\allowbreak{}services/\allowbreak{}serviceManager.js}{227--260}; (9) run a one-off repair of a
legacy \texttt{valueSat: NaN} bug \srccode{flux/\allowbreak{}ZelBack/\allowbreak{}src/\allowbreak{}services/\allowbreak{}serviceManager.js}{265--267};
(10) schedule volume-mount validation at 45\,s and hardware-requirement validation (which removes
non-compliant apps) at 50\,s, explicitly "before boot reconciliation"
\srccode{flux/\allowbreak{}ZelBack/\allowbreak{}src/\allowbreak{}services/\allowbreak{}serviceManager.js}{285--297}; (11) migrate any container still
on Docker's native \texttt{unless-stopped}\slash\texttt{always} restart policy to \texttt{no}, "because
\FluxOS manages container startup after dbReady" \srcdoc{flux/ZelBack/src/services/serviceManager.js:301}
\srccode{flux/\allowbreak{}ZelBack/\allowbreak{}src/\allowbreak{}services/\allowbreak{}serviceManager.js}{299--302}; (12) start the container
reconciler and its Docker die-event bridge, which begin queuing triggers but do not yet drain them
\srccode{flux/\allowbreak{}ZelBack/\allowbreak{}src/\allowbreak{}services/\allowbreak{}serviceManager.js}{304--310}; (13) read boot context and, fire-and-forget,
reconcile apps on boot \srccode{flux/\allowbreak{}ZelBack/\allowbreak{}src/\allowbreak{}services/\allowbreak{}serviceManager.js}{312--319}; (14) build
the \fluxd client and block on \texttt{wait\allowbreak{}For\allowbreak{}Daemon\allowbreak{}Rpc()} — the code comment notes this daemon
must be running before \texttt{daemonReady} is set, and that boot reconciliation separately
enforces an internal 5-minute timeout on the same event
\srccode{flux/\allowbreak{}ZelBack/\allowbreak{}src/\allowbreak{}services/\allowbreak{}serviceManager.js}{321--328}; (15)--(19) wire the app-sync
orchestrator, adjust the firewall, remove any stray watchtower container, and start P2P discovery
unless \texttt{config.\allowbreak{}fluxapps.\allowbreak{}discoveryAutostart === false}
\srccode{flux/\allowbreak{}ZelBack/\allowbreak{}src/\allowbreak{}services/\allowbreak{}serviceManager.js}{319--414}; (20) register the long tail of
periodic monitors listed in Table~\ref{tab:monitors}. On any thrown exception anywhere in this function,
the entire function is retried after a flat 15{,}000\,ms — not an exponential backoff — with no
retry cap \srccode{flux/\allowbreak{}ZelBack/\allowbreak{}src/\allowbreak{}services/\allowbreak{}serviceManager.js}{568--572}; boot is therefore a
self-healing infinite retry loop rather than a circuit breaker, a property revisited in
\S\ref{sec:fluxos-trust}.

\begin{table}[H]
\centering
\caption{Representative periodic monitors registered at the end of boot.}
\label{tab:monitors}
\begin{tabular}{l l}
\toprule
Monitor & Interval \\
\midrule
Firewall/UFW re-check + mount test & 30\,s \\
Stop-non-\Flux-container sweep, app-availability monitor, port restore & 1\,min (port restore repeats every 600{,}000\,ms) \\
Enterprise identity resolution & every 5\,min until resolved \\
Node-status monitor & 1.5\,min \\
CPU-usage check (\texttt{cpu\allowbreak{}Check\allowbreak{}Interval\allowbreak{}Ms}) & 900{,}000\,ms (15\,min) \\
App-availability check & 3\,min \\
Image compliance sweep & 3{,}600{,}000\,ms (1\,h) \\
Forced-removal sweep & 30\,min, then every 2\,h \\
Daemon health monitor & 5\,min, then every 15\,min \\
Storage-space check & 20\,min \\
Local-backup cleanup & 25\,min \\
\bottomrule
\end{tabular}
\srccode{flux/\allowbreak{}ZelBack/\allowbreak{}src/\allowbreak{}services/\allowbreak{}serviceManager.js}{433--563}
\end{table}

\subsubsection{Databases and collections}
\FluxOS keeps five Mongo databases, all on the local instance bound above
\srccode{flux/\allowbreak{}ZelBack/\allowbreak{}config/\allowbreak{}default.js}{33--87}:
\begin{itemize}
\item \texttt{zelfluxlocal} — session and local-node bookkeeping: \texttt{loggedusers},
  \texttt{activeloginphrases}, \texttt{activesignatures}, \texttt{activepaymentrequests},
  \texttt{completedpayments}, \texttt{geolocation}, \texttt{benchmark}, \texttt{apptamperingevents},
  \texttt{nodestartuptracker}.
\item \texttt{zelcashdata} — a mirror of \fluxd chain-scan state: \texttt{scannedheight},
  \texttt{utxoindex}, \texttt{addresstransactionindex}, \texttt{zelnodetransactions},
  \texttt{zelappshashes}, \texttt{coinbasefusionindex}.
\item \texttt{localzelapps} — this node's own truth about its containers:
  \texttt{zelappsinformation} and \texttt{zelappsruntimestate} (the per-component controller state
  consumed by the reconciler in \S\ref{sec:fluxos-reconciler}: desired state, restart
  history/crash-backoff position, last exit).
\item \texttt{globalzelapps} — network-global truth, replicated by P2P gossip and not by consensus
  (\S\ref{sec:fluxos-p2p}): \texttt{zelappsmessages}, \texttt{zelappsinformation},
  \texttt{zelappstemporarymessages}, \texttt{zelappslocation}, \texttt{appsinstallinglocations},
  \texttt{apps\allowbreak{}Installing\allowbreak{}Errors\allowbreak{}Locations}, \texttt{appstateevents} (an event log of
  \texttt{apprunning}\slash\texttt{sigterm}\slash\texttt{appremoved}\slash\texttt{evicted} transitions),
  \texttt{fluxappinstallingbroadcasts}, \texttt{fluxappinstallingerrorsbroadcasts}.
\item \texttt{chainparams} — \texttt{chainmessages}: soft-fork price/parameter messages posted
  on-chain by the Foundation multisig, active from their occurrence height (\S\ref{sec:fluxos-pricing}).
\end{itemize}

\subsubsection{HTTP API surface and the privilege model}
The route table (\texttt{ZelBack/\allowbreak{}src/\allowbreak{}routes.\allowbreak{}js}) is organized into groups, each guarded by
\texttt{verificationHelper.\allowbreak{}verifyPrivilege(privilege, req, appName)}. Every check is keyed off a
\texttt{zelidauth} header \texttt{\{zelid,\allowbreak{} signature,\allowbreak{} loginPhrase\}} verified with
\texttt{bitcoinMessage.\allowbreak{}verify} against the ZelID, a base58 P2PKH-style address
\srccode{flux/\allowbreak{}ZelBack/\allowbreak{}src/\allowbreak{}services/\allowbreak{}verificationHelper.js}{17--48,90--106}.

\begin{table}[H]
\centering
\caption{Privilege tiers enforced by \texttt{verify\allowbreak{}Privilege}.}
\begin{tabular}{p{3.15cm} p{4.95cm} p{2.2cm} p{3cm}}
\toprule
Privilege & Definition & Freshness window & Representative endpoints \\
\midrule
\texttt{admin} & \texttt{zelid === userconfig.\allowbreak{}initial.\allowbreak{}zelid} — the node's own operator
  \srccode{flux/\allowbreak{}ZelBack/\allowbreak{}src/\allowbreak{}services/\allowbreak{}verificationHelperUtils.js}{19--46} & re-signed each call & \texttt{/daemon/stop}, \texttt{/\allowbreak{}daemon/\allowbreak{}reindex},
  \texttt{/\allowbreak{}daemon/\allowbreak{}createfluxnodekey}, \texttt{/\allowbreak{}daemon/\allowbreak{}signrawtransaction} \\
\texttt{fluxteam} & \texttt{zelid} $\in \{$\texttt{fluxTeamFluxID}, \texttt{flux\allowbreak{}Support\allowbreak{}Team\allowbreak{}Flux\allowbreak{}ID}$\}$
  \srccode{flux/\allowbreak{}ZelBack/\allowbreak{}src/\allowbreak{}services/\allowbreak{}verificationHelperUtils.js}{96--122} & re-signed each call & cross-node views spanning more than this node's own data \\
\texttt{adminandfluxteam} & union of the two above
  \srccode{flux/\allowbreak{}ZelBack/\allowbreak{}src/\allowbreak{}services/\allowbreak{}verificationHelperUtils.js}{130--156} & re-signed each call & \texttt{/daemon/start}, \texttt{/\allowbreak{}daemon/\allowbreak{}restart},
  \texttt{/daemon/ping}, \texttt{/\allowbreak{}daemon/\allowbreak{}startbenchmark}, \texttt{/\allowbreak{}flux/\allowbreak{}broadcastmessage*} \\
\texttt{appowner}/ \texttt{appownerabove} & caller matches the app's current owner as resolved by
  \texttt{registryManager.\allowbreak{}getApplicationOwner(appName)}; the \texttt{-above} variant also admits
  admin/fluxteam \srccode{flux/\allowbreak{}ZelBack/\allowbreak{}src/\allowbreak{}services/\allowbreak{}verificationHelperUtils.js}{164--248} & 2\,h & \texttt{/apps/appstart}, \texttt{/apps/appstop},
  \texttt{/\allowbreak{}apps/\allowbreak{}apprestart}, \texttt{/\allowbreak{}apps/\allowbreak{}appremove}, \texttt{/apps/redeploy},
  \texttt{/\allowbreak{}apps/\allowbreak{}redeploycomponent} \\
\texttt{user} & any ZelID with a valid signature over a \texttt{loginPhrase} carrying a 13-digit
  millisecond timestamp no older than 16\,h, unless already in \texttt{loggedusers}
  \srccode{flux/\allowbreak{}ZelBack/\allowbreak{}src/\allowbreak{}services/\allowbreak{}verificationHelperUtils.js}{54--88} & 16\,h & \texttt{/\allowbreak{}apps/\allowbreak{}appregister},
  \texttt{/\allowbreak{}apps/\allowbreak{}appupdate} (self), \texttt{/\allowbreak{}daemon/\allowbreak{}prioritisetransaction} \\
\bottomrule
\end{tabular}
\end{table}

All six checks re-verify the ECDSA signature over the login phrase even when a \texttt{loggedusers}
session document exists; the session document only relaxes the timestamp-freshness check, never the
signature check \srccode{flux/\allowbreak{}ZelBack/\allowbreak{}src/\allowbreak{}services/\allowbreak{}verificationHelperUtils.js}{54--248}. The route
groups themselves are: \texttt{/daemon/*} (a thin proxy to \fluxd's JSON-RPC surface, split across
public/cached, \texttt{user}, \texttt{admin}, and \texttt{adminandfluxteam} endpoints)
\srccode{flux/\allowbreak{}ZelBack/\allowbreak{}src/\allowbreak{}routes.js}{78--1408}; \texttt{/apps/*} (public/cached reads,
\texttt{appownerabove} lifecycle actions, a \texttt{user}-level \texttt{/\allowbreak{}apps/\allowbreak{}appregister}, and an
\texttt{/\allowbreak{}apps/\allowbreak{}appupdate} reachable by the owner or by \texttt{adminandfluxteam} on the owner's
behalf) \srccode{flux/\allowbreak{}ZelBack/\allowbreak{}src/\allowbreak{}routes.js}{372--1388}; \texttt{/id/*} (ZelID session
management) \srccode{flux/\allowbreak{}ZelBack/\allowbreak{}src/\allowbreak{}routes.js}{508--518}; \texttt{/benchmark/*} (a proxy to
\fluxbench RPC) \srccode{flux/\allowbreak{}ZelBack/\allowbreak{}src/\allowbreak{}services/\allowbreak{}benchmarkService.js}{126--427}; and
\texttt{/flux/*} (network-control endpoints, including the broadcast primitives used by
\S\ref{sec:fluxos-p2p}) \srccode{flux/\allowbreak{}ZelBack/\allowbreak{}src/\allowbreak{}routes.js}{1022}.

\subsection{The application lifecycle}
\label{sec:fluxos-lifecycle}

An application moves through eight named stages between a user's signature and a running
container. Table~\ref{tab:lifecycle-timing} collects the timing constants; the paragraphs below name
each stage and its transition.

\textbf{Construct and submit.} The client signs \texttt{\{appSpecification,\allowbreak{} timestamp,\allowbreak{} signature,\allowbreak{}
type,\allowbreak{} version\}}, where \texttt{type} must be \texttt{zelappregister} or \texttt{fluxappregister}
and the envelope \texttt{version} (distinct from \texttt{appSpecification.\allowbreak{}version}) must be 1
\srccode{flux/\allowbreak{}ZelBack/\allowbreak{}src/\allowbreak{}services/\allowbreak{}appDatabase/\allowbreak{}registryManager.js}{1755--1764}, and \texttt{POST}s
it to \texttt{/\allowbreak{}apps/\allowbreak{}appregister} at \texttt{user} privilege
\srccode{flux/\allowbreak{}ZelBack/\allowbreak{}src/\allowbreak{}services/\allowbreak{}appDatabase/\allowbreak{}registryManager.js}{1714--1737}.

\textbf{Verify.} The node rejects the request if it has fewer than 8 outbound or 4 inbound peers
("not enough peers for safe application registration")
\srccode{flux/\allowbreak{}ZelBack/\allowbreak{}src/\allowbreak{}services/\allowbreak{}appDatabase/\allowbreak{}registryManager.js}{1744--1750}; rejects a
\specv{9} spec unless \texttt{config.\allowbreak{}fluxapps.\allowbreak{}acceptV9} is set (default \texttt{false}; \S10
covers the version-gating scheme in full) --- a check that exists only on an unpublished working
branch, not in the deployed source (\inprogress)
\srcbranch{RunOnFlux/\allowbreak{}flux}{feat/appspec-v9}{ZelBack/src/services/appDatabase/registryManager.js:1766-1769};
rejects a timestamp more than 1\,h old or more than 5\,min in the future
\srccode{flux/\allowbreak{}ZelBack/\allowbreak{}src/\allowbreak{}services/\allowbreak{}appDatabase/\allowbreak{}registryManager.js}{1777--1781}; and, once the
daemon is confirmed synced, runs the full schema/version gate
(\texttt{appValidator.\allowbreak{}verifyAppSpecifications}; \S10 details the version-gating rules)
\srccode{flux/\allowbreak{}ZelBack/\allowbreak{}src/\allowbreak{}services/\allowbreak{}appDatabase/\allowbreak{}registryManager.js}{1783--1798}. For enterprise
apps the signature is verified against the \emph{unformatted} spec — the signature covers the
pre-encryption content — after which \texttt{contacts}\slash\texttt{compose} are zeroed out of the
object that is actually hashed and stored
\srccode{flux/\allowbreak{}ZelBack/\allowbreak{}src/\allowbreak{}services/\allowbreak{}appDatabase/\allowbreak{}registryManager.js}{1809--1824}. The node then
computes \texttt{messageHASH} = SHA-256(type $\Vert$ version $\Vert$ JSON(spec) $\Vert$ timestamp
$\Vert$ signature) \srccode{flux/\allowbreak{}ZelBack/\allowbreak{}src/\allowbreak{}services/\allowbreak{}appDatabase/\allowbreak{}registryManager.js}{1827--1829}.

\textbf{Broadcast (temporary).} The message is stored with a 1\,h TTL and broadcast to all peers
\srccode{flux/\allowbreak{}ZelBack/\allowbreak{}src/\allowbreak{}services/\allowbreak{}appDatabase/\allowbreak{}registryManager.js}{1831--1841}. The API call then
blocks through a fixed confirmation round-trip — \texttt{delay(1200)}, request the message back
from peers, \texttt{delay(1200)}, then poll for its own existence up to 20 times at 500\,ms — and
throws \emph{"Unable to register application on the network. Try again later."} if it is still
absent \srccode{flux/\allowbreak{}ZelBack/\allowbreak{}src/\allowbreak{}services/\allowbreak{}appDatabase/\allowbreak{}registryManager.js}{1843--1858}. On success
the API returns only the message hash: this endpoint never touches the chain, and the caller must
still pay (\S\ref{sec:fluxos-pricing}) for the registration to ever become permanent
\srccode{flux/\allowbreak{}ZelBack/\allowbreak{}src/\allowbreak{}services/\allowbreak{}appDatabase/\allowbreak{}registryManager.js}{1737}.

\textbf{On-chain confirmation.} \fluxd's chain scan (\texttt{explorerService.\allowbreak{}processInsight}, run
per block) recognizes a payment by its output address — \texttt{fluxapps.\allowbreak{}address} at any height, or
one of two height-gated multisig addresses — sums \texttt{valueSat} across matching outputs, and
extracts the temp-message hash from the transaction's \texttt{OP\_RETURN}
\srccode{flux/\allowbreak{}ZelBack/\allowbreak{}src/\allowbreak{}services/\allowbreak{}explorerService.js}{309--345}, queuing
\texttt{messageVerifier.\allowbreak{}checkAndRequestApp} \srccode{flux/\allowbreak{}ZelBack/\allowbreak{}src/\allowbreak{}services/\allowbreak{}explorerService.js}{444}.

\textbf{Promote to permanent global state.} \texttt{check\allowbreak{}And\allowbreak{}Request\allowbreak{}App} re-verifies update
signatures against the current permanent owner immediately before promotion — an explicit
anti-race measure against two updates racing the same base state
\srccode{flux/\allowbreak{}ZelBack/\allowbreak{}src/\allowbreak{}services/\allowbreak{}appMessaging/\allowbreak{}messageVerifier.js}{636--682} — computes the
PoN-fork-aware expiration height, and for a fresh registration requires
\texttt{valueSat} $\geq$ \texttt{appPrice}$\times 10^{8}$. \textbf{If underpaid, the registration is
silently dropped}: the code logs a warning and returns, with no error surfaced to the payer, no
retry, and no refund path visible in this repository
\srccode{flux/\allowbreak{}ZelBack/\allowbreak{}src/\allowbreak{}services/\allowbreak{}appMessaging/\allowbreak{}messageVerifier.js}{735--763}; updates have the
same silent-drop behavior on underpayment
\srccode{flux/\allowbreak{}ZelBack/\allowbreak{}src/\allowbreak{}services/\allowbreak{}appMessaging/\allowbreak{}messageVerifier.js}{797--828}. This is a
correctness gap, not a documented design choice, and \S16 takes up what a payer-facing signal
should look like. On success, \texttt{insert\allowbreak{}App\allowbreak{}Specifications} makes the spec part of
\texttt{globalzelapps}, and any updates queued while the registration was outstanding are flushed
\srccode{flux/\allowbreak{}ZelBack/\allowbreak{}src/\allowbreak{}services/\allowbreak{}appMessaging/\allowbreak{}messageVerifier.js}{735--763}.

\textbf{Placement.} Every node independently decides whether to install (\S\ref{sec:fluxos-placement}).

\textbf{Docker deployment and running.} \texttt{appInstaller.\allowbreak{}registerAppLocally} performs port
setup, image pull/verification, and Docker network creation before calling
\texttt{install\allowbreak{}Application\allowbreak{}Hard} or \texttt{install\allowbreak{}Application\allowbreak{}Soft}
\srccode{flux/\allowbreak{}ZelBack/\allowbreak{}src/\allowbreak{}services/\allowbreak{}appLifecycle/\allowbreak{}appInstaller.js}{415--746}; the evidence available
does not trace the exact Docker API calls (container-create options, network-attach order, volume
mount construction) inside those two functions beyond their call sites
\srcinferred. Container start/stop itself is centralized in the reconciler
(\S\ref{sec:fluxos-reconciler}), not in the installer.

\textbf{Monitoring and expiry.} While running, the app is subject to the node-health, benchmark,
and image-compliance gates of \S\ref{sec:fluxos-resources} and \S\ref{sec:fluxos-moderation}, and —
where the operator has opted the app into it — per-container telemetry shipped by
\texttt{flux-\allowbreak{}telemetryd} to a customer-chosen backend, sampled every 15\,s
\srccode{flux-telemetryd/\allowbreak{}src/\allowbreak{}identity.rs}{1--31,56--74}
\srccode{flux-telemetryd/\allowbreak{}src/\allowbreak{}collector.rs}{26} \srcdoc{flux-telemetryd/README.md:1-20} (\shipped, a
companion daemon evidenced among the node-level services surrounding \FluxOS rather than in
\FluxOS itself). A departing node's \texttt{fluxnodesigterm} broadcast extends its apps' location
TTL by \texttt{SIGTERM\_EXPIRY\_MS} = 420{,}000\,ms (7\,min) rather than expiring them immediately
(\S\ref{sec:fluxos-p2p}); a companion node-wide/per-app graceful-drain daemon,
\texttt{flux-shutdownd}, exists to broadcast a machine-readable termination reason to the container
before this happens, but only its signal-and-cleanup pipeline stages are wired to production
events today — drain-broadcast and \texttt{preStop} execution are modelled and unit-tested but not
yet reachable \srcdoc{flux-shutdownd/crates/shutdownd/src/pipeline/state\_machine.rs:3-7}
(\inprogress). An app's actual expiration height, once reached, is
enforced by a periodic expiry pass every \texttt{expire\allowbreak{}Flux\allowbreak{}Apps\allowbreak{}Period} = 100 blocks
\srccode{flux/\allowbreak{}ZelBack/\allowbreak{}config/\allowbreak{}default.js}{298}.

\begin{table}[H]
\centering
\caption{Lifecycle timing constants.}
\label{tab:lifecycle-timing}
\begin{tabular}{l l}
\toprule
Stage / constant & Value \\
\midrule
Registration timestamp window & $>1$\,h old or $>5$\,min future rejected \srccode{flux/\allowbreak{}ZelBack/\allowbreak{}src/\allowbreak{}services/\allowbreak{}appDatabase/\allowbreak{}registryManager.js}{1777--1781} \\
Temp-message TTL & 3{,}600\,s \srccode{flux/\allowbreak{}ZelBack/\allowbreak{}config/\allowbreak{}default.js}{314} \\
Registration confirmation round-trip & $2\times1{,}200$\,ms delay $+$ up to $20\times500$\,ms poll ($\approx$8--10\,s) \srccode{flux/\allowbreak{}ZelBack/\allowbreak{}src/\allowbreak{}services/\allowbreak{}appDatabase/\allowbreak{}registryManager.js}{1843--1858} \\
\texttt{minOutgoing} / \texttt{minIncoming} peers to register & 8 / 4 \srccode{flux/\allowbreak{}ZelBack/\allowbreak{}config/\allowbreak{}default.js}{262--265} \\
Location TTL & 7{,}500\,s \srccode{flux/\allowbreak{}ZelBack/\allowbreak{}config/\allowbreak{}default.js}{311} \\
Installing-broadcast TTL & 900\,s \srccode{flux/\allowbreak{}ZelBack/\allowbreak{}config/\allowbreak{}default.js}{312} \\
Install-error TTL & 3{,}600\,s \srccode{flux/\allowbreak{}ZelBack/\allowbreak{}config/\allowbreak{}default.js}{313} \\
SIGTERM grace window & 420{,}000\,ms (7\,min) \srccode{flux/\allowbreak{}ZelBack/\allowbreak{}src/\allowbreak{}services/\allowbreak{}utils/\allowbreak{}appConstants.js}{86} \\
\texttt{blocksLasting} (default lifetime) & 22{,}000 blocks ($\approx$1 month pre-\PoN, $\approx$1 week post-\PoN) \srccode{flux/\allowbreak{}ZelBack/\allowbreak{}config/\allowbreak{}default.js}{285} \\
Expiry sweep period & 100 blocks \srccode{flux/\allowbreak{}ZelBack/\allowbreak{}config/\allowbreak{}default.js}{298} \\
Update sweep / removal sweep period & 9 / 11 blocks \srccode{flux/\allowbreak{}ZelBack/\allowbreak{}config/\allowbreak{}default.js}{299--300} \\
\bottomrule
\end{tabular}
\end{table}

\subsection{Placement: opportunistic self-selection with broadcast collision-avoidance}
\label{sec:fluxos-placement}

There is no hash-seeded, deterministic node-selection algorithm, and no component that assigns an
application to a node. Every node runs the same self-looping procedure
(\texttt{appSpawner.\allowbreak{}trySpawningGlobalApplication}, invoked as \texttt{spawnLoop} with the delay
returned by its previous iteration) \srccode{flux/\allowbreak{}ZelBack/\allowbreak{}src/\allowbreak{}services/\allowbreak{}appLifecycle/\allowbreak{}appSpawner.js}{54--65,77--787},
against its own local read of the gossiped \texttt{globalzelapps} state. This is the mechanism the
outline's single claim is about, and Algorithm~\ref{alg:placement} states it in full.

\begin{definition}[Opportunistic self-selection]
A node that observes an application with fewer running instances than its specification requests
may attempt to become one of those instances. It decides to attempt \emph{unilaterally}: no other
node, and no on-chain transaction, grants or denies it that attempt in advance. Conflicts among
simultaneous attempts are resolved after the fact, by broadcast and a wait, not before the fact by
assignment.
\end{definition}

\begin{algorithm}[H]
\footnotesize
\DontPrintSemicolon
\caption{One iteration of the per-node placement loop
  \srccode{flux/\allowbreak{}ZelBack/\allowbreak{}src/\allowbreak{}services/\allowbreak{}appLifecycle/\allowbreak{}appSpawner.js}{77--787}}
\label{alg:placement}
\KwData{this node's local state; its local replica of \texttt{globalzelapps}}
\If{\textbf{not} (enterprise identity resolved \textbf{and} daemon synced \textbf{and} local DB
    ready \textbf{and} \textbf{not} DoS-flagged \textbf{and} node confirmed \textbf{and}
    \fluxbench not erroring \textbf{and} own external IP known \textbf{and} locally installed
    apps $<$ \texttt{maxAppsPerNode} = 200)}{
  \Return (retry after the loop's error-path delay) \tcp*{gate, \srccode{flux/\allowbreak{}ZelBack/\allowbreak{}src/\allowbreak{}services/\allowbreak{}appLifecycle/\allowbreak{}appSpawner.js}{78--150}}
}
$M \gets$ apps whose actual (\PoN-fork-aware) expiration height has not passed and whose running
  instance count is below \texttt{instances} $\lor 3$ \tcp*{Mongo aggregation, \srccode{flux/\allowbreak{}ZelBack/\allowbreak{}src/\allowbreak{}services/\allowbreak{}appLifecycle/\allowbreak{}appSpawner.js}{152--247}}
\eIf{a due entry exists in the deferred-recheck queue \texttt{apps\allowbreak{}To\allowbreak{}Be\allowbreak{}Checked\allowbreak{}Later}}{
  $\app \gets$ pop that entry \tcp*{deterministic, not random}
}{
  $C \gets \{a \in M : a.\texttt{nodes} \text{ names this node's own IP}\}$\;
  \eIf{$C \neq \emptyset$}{$\app \gets$ uniform-random draw from $C$}{$\app \gets$ uniform-random
    draw from $M$} \tcp*{\srccode{flux/\allowbreak{}ZelBack/\allowbreak{}src/\allowbreak{}services/\allowbreak{}appLifecycle/\allowbreak{}appSpawner.js}{254--336}: own-IP bias, then uniform}
}
\If{$\app$ fails a node-targeting or geolocation filter}{
  discard $\app$ this cycle, continue the loop
  \tcp*{enterprise \texttt{nodes} is strict at any version; non-enterprise \texttt{nodes} binds
  only versions $<8$; geolocation allow/deny by \texttt{ac<geo>}\slash\texttt{a!c<geo>} prefix match
  (exact encoding not fully traced) \srccode{flux/\allowbreak{}ZelBack/\allowbreak{}src/\allowbreak{}services/\allowbreak{}appLifecycle/\allowbreak{}appSpawner.js}{288--300} \srcinferred}
}
\If{$\app$ trips a soft-deferral condition (non-enterprise app on an \ArcaneOS node; static-IP
    mismatch; datacenter mismatch; targeted-\texttt{nodes} app not naming this
    node; a capacity-gap tier where this node is too large for a small app)}{
  push $\app$ onto \texttt{apps\allowbreak{}To\allowbreak{}Be\allowbreak{}Checked\allowbreak{}Later} with a condition- and enterprise-specific delay
  (Table~\ref{tab:spawn-deferrals}); \Return \tcp*{\srccode{flux/\allowbreak{}ZelBack/\allowbreak{}src/\allowbreak{}services/\allowbreak{}appLifecycle/\allowbreak{}appSpawner.js}{520--630}}
}
verify the image against the Docker Hub whitelist and pull it; verify port availability, including
  a public reachability probe relayed through peers
  \tcp*{\texttt{portManager.\allowbreak{}checkInstallingAppPortAvailable}}
broadcast \texttt{fluxappinstalling} $\{$type, version:\,1, name, own IP, \texttt{broadcastedAt}$\}$
  to all peers and store it locally\;
wait \texttt{install\allowbreak{}Collision\allowbreak{}Wait\allowbreak{}Ms} = 90{,}000\,ms \tcp*{\srccode{flux/\allowbreak{}ZelBack/\allowbreak{}config/\allowbreak{}default.js}{321}, comment: "give it 1.5m so messages are propagated"}
re-read the global installing $+$ running location lists for $\app$\;
\If{combined count $>$ required instances}{
  rank all installing nodes by \texttt{broadcastedAt} ascending\;
  \If{this node's rank does not fit within the remaining slots}{abandon $\app$; \Return}
  \tcp*{\srccode{flux/\allowbreak{}ZelBack/\allowbreak{}src/\allowbreak{}services/\allowbreak{}appLifecycle/\allowbreak{}appSpawner.js}{686--703}: earliest-broadcast wins — the entire conflict-resolution mechanism}
}
$\mathit{ok} \gets$ \texttt{appInstaller.\allowbreak{}registerAppLocally}($\app$) \tcp*{Docker deployment}
\If{\textbf{not} $\mathit{ok}$}{
  cache $\app$'s hash in \texttt{spawn\allowbreak{}Errors\allowbreak{}Longer\allowbreak{}App\allowbreak{}Cache} (no automatic retry until it expires);
  \Return \tcp*{\srccode{flux/\allowbreak{}ZelBack/\allowbreak{}src/\allowbreak{}services/\allowbreak{}appLifecycle/\allowbreak{}appSpawner.js}{730--743}}
}
wait 60\,s; re-read running locations for $\app$; rank by \texttt{runningSince} (nodes with none
  yet sort first)\;
\If{this node's rank $>$ \texttt{minInstances}}{
  uninstall $\app$ locally, immediately \tcp*{self-correction, \srccode{flux/\allowbreak{}ZelBack/\allowbreak{}src/\allowbreak{}services/\allowbreak{}appLifecycle/\allowbreak{}appSpawner.js}{745--775}}
}
\end{algorithm}

\begin{table}[H]
\centering
\caption{Soft-deferral delays (push back into \texttt{apps\allowbreak{}To\allowbreak{}Be\allowbreak{}Checked\allowbreak{}Later}, do not abandon).}
\label{tab:spawn-deferrals}
\begin{tabular}{l l l}
\toprule
Condition & Enterprise app & Standard app \\
\midrule
Non-enterprise app on an \ArcaneOS node & — & 120{,}000\,ms \srccode{flux/\allowbreak{}ZelBack/\allowbreak{}src/\allowbreak{}services/\allowbreak{}appLifecycle/\allowbreak{}appSpawner.js}{34} \\
Targeted \texttt{nodes} not naming this node & 1{,}800{,}000\,ms & 3{,}420{,}000\,ms \srccode{flux/\allowbreak{}ZelBack/\allowbreak{}config/\allowbreak{}default.js}{341} \\
Static-IP mismatch & 1{,}620{,}000\,ms & 3{,}420{,}000\,ms \srccode{flux/\allowbreak{}ZelBack/\allowbreak{}config/\allowbreak{}default.js}{342} \\
Datacenter mismatch & 1{,}620{,}000\,ms & 3{,}420{,}000\,ms \srccode{flux/\allowbreak{}ZelBack/\allowbreak{}config/\allowbreak{}default.js}{343} \\
Capacity gap, large/medium/small & 1{,}800{,}000 / 1{,}260{,}000 / 720{,}000\,ms & 7{,}020{,}000 / 5{,}220{,}000 / 3{,}420{,}000\,ms \srccode{flux/\allowbreak{}ZelBack/\allowbreak{}config/\allowbreak{}default.js}{345--347} \\
\bottomrule
\end{tabular}
\end{table}

The datacenter-mismatch row exists in the deferral table but the branch that would trigger it is
flagged by the code's own comment as currently unreachable: apps with \texttt{datacenter=\allowbreak{}true} are
routed through an ownership filter to enterprise nodes before this deferral chain runs, so
\texttt{datacenter} is always false by the time this check executes
\srcdoc{flux/ZelBack/src/services/appLifecycle/appSpawner.js:568}. A syncthing-specific tie-break
additionally picks the earliest broadcaster among nodes sharing a \texttt{/24}-like IP-prefix range
when several attempt installation concurrently, so replicated storage does not converge two masters
on the same subnet \srccode{flux/\allowbreak{}ZelBack/\allowbreak{}src/\allowbreak{}services/\allowbreak{}appLifecycle/\allowbreak{}appSpawner.js}{706--727}.

\begin{property}[Earliest-broadcast-wins is the entire conflict-resolution mechanism]
When more nodes attempt installation than an application's specification requests, the network
resolves the conflict by ranking the competing \texttt{fluxappinstalling} broadcasts by their
self-reported \texttt{broadcastedAt} timestamp and letting only the earliest ones proceed. There is
no leader election, no quorum, and no on-chain arbitration anywhere in this path
\srccode{flux/\allowbreak{}ZelBack/\allowbreak{}src/\allowbreak{}services/\allowbreak{}appLifecycle/\allowbreak{}appSpawner.js}{686--703}.
\end{property}

\textbf{Invariant and failure mode.} The mechanism's target property is that, given enough peer
connectivity for gossip to be reliable (the spawner is only active above an
\texttt{app\allowbreak{}Sync\allowbreak{}Peer\allowbreak{}Threshold} of 12 live peers and is explicitly paused below an
\texttt{app\allowbreak{}Sync\allowbreak{}Degraded\allowbreak{}Threshold} of 4, "gossip unreliable"
\srccode{flux/\allowbreak{}ZelBack/\allowbreak{}config/\allowbreak{}default.js}{268--269}, \S\ref{sec:fluxos-p2p}), the number of nodes
that both install and keep an application converges to its requested instance count within roughly
one 90-second wait plus one 60-second post-install check. The mechanism's honesty assumption is
exactly as narrow as that sentence: it depends on every competing node broadcasting an honest
\texttt{broadcastedAt} and actually backing off when its rank does not fit. Nothing in the code
paths read re-verifies \texttt{broadcastedAt} against the signed message envelope's own timestamp
for this specific message type \srcinferred; a node that broadcasts a forged early
\texttt{broadcastedAt}, or that never re-checks its rank, can over-claim install slots, and a
network partition that delays propagation past the fixed 90{,}000\,ms window can produce the same
effect without any dishonesty at all. There is no mechanism above this one — no on-chain record, no
elected coordinator — that would notice or correct such an outcome directly; the only correction is
the pair of independent loops described next.

\begin{remark}[Convergence by opposing correction, not by agreement]
The system does not converge on the target instance count because nodes agree on who should run an
application. It converges because two independent, opposing loops keep nudging the count back
toward the target from either side. First, immediately after its own install, a node re-checks its
rank by \texttt{runningSince} (nodes with no \texttt{runningSince} yet — i.e. still starting —
sort first) and uninstalls itself at once if it now ranks below \texttt{minInstances}
\srccode{flux/\allowbreak{}ZelBack/\allowbreak{}src/\allowbreak{}services/\allowbreak{}appLifecycle/\allowbreak{}appSpawner.js}{745--775}. Second, independently and
periodically, \texttt{advancedWorkflows.\allowbreak{}checkAndRemoveApplicationInstance} scans every locally
installed app whose live location count exceeds \texttt{instances}, ranks all running instances by
\texttt{runningSince} descending (newest first) with IP descending as a tiebreak, and removes
locally if this node ranks 0 — the newest — waiting \texttt{removal.delay} = 300\,s before
evaluating the next over-provisioned app "so we don't have more removals at the same time"
\srccode{flux/\allowbreak{}ZelBack/\allowbreak{}src/\allowbreak{}services/\allowbreak{}appLifecycle/\allowbreak{}advancedWorkflows.js}{2899--2966,2928--2963}. Both
loops correct over-installation; neither one, nor anything else in this design, corrects
under-installation faster than the next node's own random draw happens to land on the
under-replicated app. Under a network partition that produces two disjoint views of
\texttt{runningSince} ordering, both loops can independently conclude "I am over-ranked," causing
more removals than intended, or conclude the opposite and remove none, until gossip re-converges.
\end{remark}

\begin{figure}[H]
\centering
\begin{tikzpicture}[
  font=\scriptsize, node distance=4.2mm,
  stp/.style={draw, rounded corners=2pt, align=center, text width=52mm, minimum height=6.5mm,
              inner ysep=2pt},
  cross/.style={stp, very thick, fill=black!8},
  outbox/.style={draw, rounded corners=2pt, align=center, text width=24mm, fill=black!4,
              minimum height=6mm, font=\scriptsize},
  arr/.style={-{Latex[length=1.8mm]}, thick},
  lbl/.style={font=\tiny, inner sep=1.5pt, align=left},
]
  \node[stp]                     (gate)  {precondition gate\\{\tiny synced, DB ready, not in DOS, confirmed, benchmark ok, under \texttt{maxAppsPerNode}}};
  \node[stp, below=of gate]      (pick)  {candidate selection\\{\tiny own-IP bias first, then \textbf{uniform random}}};
  \node[stp, below=of pick]      (filt)  {targeting and geolocation filters};
  \node[stp, below=of filt]      (defer) {soft-deferral branches\\{\tiny capacity tiers, static IP, ArcaneOS, datacenter}};
  \node[cross, below=of defer]   (bcast) {broadcast \texttt{fluxappinstalling}\\{\tiny then wait a fixed \SI{90000}{\milli\second}}};
  \node[cross, below=of bcast]   (rank)  {re-read global list; rank by \texttt{broadcastedAt}\\{\tiny proceed only if this node still fits the remaining slots}};
  \node[stp, below=of rank]      (inst)  {install locally (Docker)};
  \node[stp, below=of inst]      (self)  {after \SI{1}{\minute}: rank by \texttt{runningSince}\\{\tiny if over-ranked, \textbf{uninstall self}}};

  \foreach \x/\y in {gate/pick, pick/filt, filt/defer, defer/bcast, bcast/rank, rank/inst, inst/self}
    \draw[arr] (\x) -- (\y);

  \node[outbox, right=7mm of defer] (dfr) {deferred: retry later};
  \node[outbox, right=7mm of rank]  (aba) {abandon this attempt};
  \draw[arr] (defer) -- (dfr);
  \draw[arr] (rank)  -- (aba);

  \draw[arr] (self.west) -- ++(-5mm,0) |- node[lbl, left, pos=0.25] {loop} (gate.west);

  \node[draw, dashed, rounded corners=2pt, inner sep=2.5mm, fit=(bcast)(rank),
        label={[font=\tiny, align=center]right:{the \emph{only} cross-node\\coupling: two broadcast\\/read points}}] {};
\end{tikzpicture}
\caption{Placement flow for one node's attempt at one application. The loop is private to a single
node: every step is a local decision, and the only coupling to other nodes is the broadcast/read pair
shown dashed. Conflicts are resolved by earliest-broadcast-wins, not by election or arbitration —
there is no scheduler in this picture, and that is the point \shipped
\srccode{flux/\allowbreak{}ZelBack/\allowbreak{}src/\allowbreak{}services/\allowbreak{}appLifecycle/\allowbreak{}appSpawner.js}{77--787}.}
\label{fig:placement-flow}
\end{figure}

\subsection{The reconciler: a single, level-based actuator}
\label{sec:fluxos-reconciler}

Once a container exists, everything that starts or stops it goes through one component,
documented in its own source as "the single, level-based actuator for app containers"
\srcdoc{flux/ZelBack/src/services/appMonitoring/appReconciler.js:20-27}. Every trigger — a Docker
\texttt{die} event, a stream reconnect, an hourly tick, boot, a post-install callback, or a
master/slave (syncthing) decision — only \emph{enqueues} a component identifier;
\texttt{reconcile()} is the only function that calls \texttt{appDockerStart}\slash\texttt{appDockerStop}
\srccode{flux/\allowbreak{}ZelBack/\allowbreak{}src/\allowbreak{}services/\allowbreak{}appMonitoring/\allowbreak{}appReconciler.js}{20--27}.

Desired state is resolved from four inputs, in strict priority order, highest wins
\srcdoc{flux/ZelBack/src/services/appMonitoring/appReconciler.js:29-45}:
\begin{enumerate}
\item \texttt{operator\allowbreak{}Stopped} — durable, persisted in \texttt{apps\allowbreak{}Runtime\allowbreak{}State}. The source calls
  this "user lock, wins over all."
\item \texttt{controller\allowbreak{}Desired} — in-memory only, re-derived every cycle by the master/slave and
  syncthing deciders. The source states why it is \emph{not} persisted: "a stale election after a
  reboot must not act."
\item \texttt{dataDesired} — an in-memory "clear" flag requesting an app-data wipe before the next
  run, also not persisted, for the same class of reason: "a stale wipe intent surviving a restart
  could delete the only good copy."
\item Restart policy and the last exit code, as the residual default.
\end{enumerate}
The syncthing/master-slave decider that feeds \texttt{controller\allowbreak{}Desired} is a second, large state
machine (\texttt{syncthingFolderStateMachine.\allowbreak{}js}, \texttt{syncthing\allowbreak{}Monitor.\allowbreak{}js}) that this evidence
pass did not read in depth; only its output channel into the reconciler is characterized here
\srcinferred.

Concurrency is single-flight per component identifier: an \texttt{inFlight} set prevents two
reconciles of the same component from running at once, and a \texttt{dirty} set forces one more
pass if the identifier was re-enqueued while a reconcile was already running
\srccode{flux/\allowbreak{}ZelBack/\allowbreak{}src/\allowbreak{}services/\allowbreak{}appMonitoring/\allowbreak{}appReconciler.js}{51--52}. A boot drain gate (an
\texttt{AsyncGate}) holds all boot-time triggers in \texttt{bootPending} until every boot-held
component has completed \emph{one} reconcile pass — not until all are running — capped at
\texttt{BOOT\_DRAIN\_SETTLE\_CAP\_MS} = 2\,min, so a single wedged reconcile can never indefinitely
suppress the node's own network-presence broadcast
\srccode{flux/\allowbreak{}ZelBack/\allowbreak{}src/\allowbreak{}services/\allowbreak{}appMonitoring/\allowbreak{}appReconciler.js}{55--64}. A container found
detached from its own Docker network — a stale libnetwork endpoint that a plain start cannot repair
— is force-removed and recreated rather than merely restarted
\srcdoc{flux/ZelBack/src/services/appMonitoring/appReconciler.js:24-26}. Crash recovery follows a
fixed backoff ladder, \texttt{crash\allowbreak{}Backoff\allowbreak{}Delays\allowbreak{}Ms} = [0, 30{,}000, 300{,}000, 900{,}000,
1{,}800{,}000]\,ms (0\,s, 30\,s, 5\,min, 15\,min, 30\,min), reset to the first rung after
\texttt{crash\allowbreak{}Backoff\allowbreak{}Stable\allowbreak{}Run\allowbreak{}Ms} = 600{,}000\,ms (10\,min) of stable running
\srccode{flux/\allowbreak{}ZelBack/\allowbreak{}config/\allowbreak{}default.js}{130--131}.

\subsection{Soft versus hard redeploy}
\label{sec:fluxos-redeploy}

Redeploy, removal, and installation are mutually exclusive with each other through four boolean
globals — \texttt{globalState.\allowbreak{}removalInProgress}, \texttt{installation\allowbreak{}In\allowbreak{}Progress},
\texttt{soft\allowbreak{}Redeploy\allowbreak{}In\allowbreak{}Progress}, \texttt{hard\allowbreak{}Redeploy\allowbreak{}In\allowbreak{}Progress} — checked at entry. If any is
already set, the call returns a warning and does nothing; it does not queue itself and does not
retry \srccode{flux/\allowbreak{}ZelBack/\allowbreak{}src/\allowbreak{}services/\allowbreak{}appLifecycle/\allowbreak{}advancedWorkflows.js}{1245--1281,1355--1383}.

\begin{table}[H]
\centering
\caption{Soft versus hard redeploy.}
\begin{tabular}{l l l}
\toprule
& Soft redeploy & Hard redeploy \\
\midrule
Teardown & \texttt{soft\allowbreak{}Remove\allowbreak{}App\allowbreak{}Locally}: stops/removes the container,
  \textbf{keeps} the app volume/data & \texttt{appUninstaller.\allowbreak{}removeAppLocally(force=false)}: full
  teardown, \textbf{drops} volumes \\
Wait & \texttt{config.\allowbreak{}fluxapps.\allowbreak{}redeploy.\allowbreak{}delay} = 30\,s & 30\,s (same constant) \\
Rebuild & \texttt{check\allowbreak{}App\allowbreak{}Requirements} then
  \texttt{soft\allowbreak{}Register\allowbreak{}App\allowbreak{}Locally} (recreate container against existing data) &
  \texttt{check\allowbreak{}App\allowbreak{}Requirements} then \texttt{appInstaller.registerAppLocally(...,\allowbreak{}
  sendRemovalMessage=true)} (full fresh install) \\
\bottomrule
\end{tabular}
\srccode{flux/\allowbreak{}ZelBack/\allowbreak{}src/\allowbreak{}services/\allowbreak{}appLifecycle/\allowbreak{}advancedWorkflows.js}{1313--1330,1394--1420}
\end{table}

\textbf{Automatic escalation.} For \specv{8}-and-later apps, if the installed component structure
differs from the target specification — component count or composition changed —
\texttt{softRedeploy} escalates itself to \texttt{hardRedeploy} instead of proceeding, "for
component structure safety" \srccode{flux/\allowbreak{}ZelBack/\allowbreak{}src/\allowbreak{}services/\allowbreak{}appLifecycle/\allowbreak{}advancedWorkflows.js}{1281--1305}.
If either path fails, the app is force-removed
(\texttt{removeAppLocally(.\allowbreak{}.\allowbreak{}.\allowbreak{}, force=true, sendMessage=true)}) as cleanup
\srccode{flux/\allowbreak{}ZelBack/\allowbreak{}src/\allowbreak{}services/\allowbreak{}appLifecycle/\allowbreak{}advancedWorkflows.js}{1332--1341,1422--1427}. Both
paths stream progress back over the original HTTP response (\texttt{res.write}) when one is
available: this is a synchronous, long-running HTTP call, not a job queue a client polls.

\subsection{The P2P layer}
\label{sec:fluxos-p2p}

\subsubsection{Wire format, dispatch, dedup}
Every \FluxOS-to-\FluxOS message is JSON over WebSocket, wrapped in a signed envelope
\texttt{\{version:\,\allowbreak{}1,\allowbreak{} timestamp,\allowbreak{} pubKey,\allowbreak{} signature (over version+message+timestamp),\allowbreak{} data\}}
\srccode{flux/\allowbreak{}ZelBack/\allowbreak{}src/\allowbreak{}services/\allowbreak{}utils/\allowbreak{}fluxBroadcastHelper.js}{11--26}. Observed
\texttt{data.type} values: \texttt{zelappregister}\slash\texttt{fluxappregister},
\texttt{zelappupdate}\slash\texttt{fluxappupdate}, \texttt{fluxapprequest} (a v1 single-hash form and a
v2 form batching up to 500 hashes), \texttt{fluxapprunning}, \texttt{fluxipchanged},
\texttt{fluxappremoved}, \texttt{fluxappinstalling}, \texttt{fluxappinstallingerror},
\texttt{fluxnodesigterm}, and four sync-round types
(\texttt{fluxapptempsync}, \texttt{fluxapprunningsync}, \texttt{fluxappinstallingsync},
\texttt{fluxappinstallingerrorssync}) \srccode{flux/\allowbreak{}ZelBack/\allowbreak{}src/\allowbreak{}services/\allowbreak{}fluxCommunication.js}{585--602,644--660}
\srccode{flux/\allowbreak{}ZelBack/\allowbreak{}src/\allowbreak{}services/\allowbreak{}fluxCommunicationMessagesSender.js}{53--74}. A separate raw,
non-JSON channel — \texttt{message\allowbreak{}Hash\allowbreak{}Present} (a 40-char hex hash) and
\texttt{request\allowbreak{}Message\allowbreak{}Hash} — bypasses the signed envelope entirely, because it only ever carries a
hash and never application data \srccode{flux/\allowbreak{}ZelBack/\allowbreak{}src/\allowbreak{}services/\allowbreak{}fluxCommunication.js}{511--536}.

Inbound dispatch (\texttt{dispatch\allowbreak{}Flux\allowbreak{}Message}) runs, per message: raw hash-gossip frames first,
handled and returned immediately, no signature check
\srccode{flux/\allowbreak{}ZelBack/\allowbreak{}src/\allowbreak{}services/\allowbreak{}fluxCommunication.js}{511--536}; a malformed envelope closes the
connection \srccode{flux/\allowbreak{}ZelBack/\allowbreak{}src/\allowbreak{}services/\allowbreak{}fluxCommunication.js}{538--546}; deduplication after
a random 1--75\,ms jitter delay against an in-memory \texttt{messageCache} keyed on
\texttt{hash(data)} — a seen hash is silently dropped
\srccode{flux/\allowbreak{}ZelBack/\allowbreak{}src/\allowbreak{}services/\allowbreak{}fluxCommunication.js}{549--554} (the exact TTL/eviction policy
of this cache was not traced to its definition in the evidence available \srcinferred); a blocklist
check against \texttt{wsPeerCache}
\srccode{flux/\allowbreak{}ZelBack/\allowbreak{}src/\allowbreak{}services/\allowbreak{}fluxCommunication.js}{556--563}; then signature/timestamp
verification, yielding \texttt{OK}, \texttt{NODE\_NOT\_FOUND} (originator not in the deterministic
node list — dropped without penalizing the relaying peer, "the relay peer is not at fault"), or a
bad-signature/malformed result, which is tracked per-peer in a rolling 10-minute window and closes
the connection at 5 or more bad messages in that window (the timestamp array is capped at 20
entries, truncated to the last 10, bounding memory from an adversarial peer)
\srccode{flux/\allowbreak{}ZelBack/\allowbreak{}src/\allowbreak{}services/\allowbreak{}fluxCommunication.js}{592--618}. A valid, fresh message is
dispatched by type via \texttt{setImmediate} — deferred one tick, not otherwise queued or
rate-limited per type \srccode{flux/\allowbreak{}ZelBack/\allowbreak{}src/\allowbreak{}services/\allowbreak{}fluxCommunication.js}{604--618}; a stale
but otherwise valid message gets a NAK rather than a silent drop
\srccode{flux/\allowbreak{}ZelBack/\allowbreak{}src/\allowbreak{}services/\allowbreak{}fluxCommunication.js}{601}. Four message types
(\texttt{fluxapprunning}, \texttt{fluxipchanged}, \texttt{fluxappremoved}, \texttt{fluxnodesigterm}) use a 240{,}000\,ms
(4\,min) freshness window \srccode{flux/\allowbreak{}ZelBack/\allowbreak{}src/\allowbreak{}services/\allowbreak{}fluxCommunication.js}{317,405,437,465}.

\subsubsection{Rebroadcast and convergence}
Rebroadcast follows one of two mutually exclusive strategies, selected by comparing daemon height
to \texttt{messages\allowbreak{}Broadcast\allowbreak{}Refactor\allowbreak{}Start} = 1{,}751{,}250 ("expected block at 13th October 2024"
per the code's own comment \srcdoc{flux/ZelBack/config/default.js:124}): below it,
\texttt{fluxCommunicationMessagesSender.\allowbreak{}relay(.\allowbreak{}.\allowbreak{}.\allowbreak{})} flood-fills the full signed message to every
other peer; at or above it, only the 40-character hash is gossiped
(\texttt{peerManager.\allowbreak{}broadcastHash}), and a receiver that lacks the body pulls it on demand via
\texttt{request\allowbreak{}Message\allowbreak{}Hash} \srccode{flux/\allowbreak{}ZelBack/\allowbreak{}src/\allowbreak{}services/\allowbreak{}fluxCommunication.js}{407--412,431--436,486--491}.
This is a live protocol transition keyed on chain height, not a node-local configuration flag. A
\texttt{fluxnodesigterm} broadcast extends the affected apps' \texttt{expireAt} in
\texttt{zelappslocation} to \texttt{broadcastedAt} $+$ \texttt{SIGTERM\_EXPIRY\_MS}
(420{,}000\,ms, 7\,min) rather than letting the location expire immediately
\srccode{flux/\allowbreak{}ZelBack/\allowbreak{}src/\allowbreak{}services/\allowbreak{}fluxCommunication.js}{472--478}
\srccode{flux/\allowbreak{}ZelBack/\allowbreak{}src/\allowbreak{}services/\allowbreak{}utils/\allowbreak{}appConstants.js}{86}. Beyond continuous gossip, a
periodic hash-sync protocol runs with its own bounded parameters — up to
\texttt{hash\allowbreak{}Sync\allowbreak{}Max\allowbreak{}Rounds} = 4 rounds, \texttt{hash\allowbreak{}Sync\allowbreak{}Peers\allowbreak{}Per\allowbreak{}Round} = 3,
\texttt{hash\allowbreak{}Sync\allowbreak{}Ephemeral\allowbreak{}Peers} = 5, \texttt{hash\allowbreak{}Sync\allowbreak{}Max\allowbreak{}Retries} = 3 at
\texttt{hash\allowbreak{}Sync\allowbreak{}Retry\allowbreak{}Ms} = 300{,}000\,ms, with a \texttt{hash\allowbreak{}Sync\allowbreak{}Settle\allowbreak{}Ms} = 4{,}000\,ms settle
window \srccode{flux/\allowbreak{}ZelBack/\allowbreak{}config/\allowbreak{}default.js}{355--363}; the exact message sequence of one full
round (who is asked, in what order, what happens on a partial response) was not traced beyond these
constants \srcinferred.

\textbf{This is gossip convergence, not consensus.} Global application state — which node is
running, installing, or has removed which application — lives in \texttt{globalzelapps}, is
propagated exclusively by the dedup/rebroadcast/hash-sync mechanism above, TTL-expires by the
constants in Table~\ref{tab:lifecycle-timing}, and is not committed to any block, checked against any
quorum, or given any finality. Two nodes with a partitioned or lagging gossip view can disagree
about the current state of an application for as long as the partition lasts; nothing above the
mesh corrects that disagreement faster than the mesh itself re-converges.

\subsubsection{Peer topology and liveness}
\texttt{Flux\allowbreak{}Peer\allowbreak{}Manager} is a single in-memory registry of inbound and outbound sockets, ephemeral
sync peers, a reconnect queue tracking per-peer attempts, an unstable-node tracker, IP-group/unique-IP
counters, and a bounded circular history buffer of peer lifecycle events
\srccode{flux/\allowbreak{}ZelBack/\allowbreak{}src/\allowbreak{}services/\allowbreak{}utils/\allowbreak{}FluxPeerManager.js}{31--77}. Inbound admission is capped at
the larger of \texttt{maxPeers} = $4\times$\texttt{minIncoming} and \texttt{max\allowbreak{}Number\allowbreak{}Of\allowbreak{}Connections}, which is normally
$9\times$\texttt{minIncoming} but scales down when \texttt{number\allowbreak{}Of\allowbreak{}Flux\allowbreak{}Nodes}/160 falls below that
figure, so small networks get a proportionally smaller inbound cap rather than a flat one, floored at \texttt{maxPeers}
\srccode{flux/\allowbreak{}ZelBack/\allowbreak{}src/\allowbreak{}services/\allowbreak{}utils/\allowbreak{}FluxPeerManager.js}{839--847}. Reconnect backoff follows
[120{,}000, 300{,}000, 600{,}000, 900{,}000]\,ms (2/5/10/15\,min)
\srccode{flux/\allowbreak{}ZelBack/\allowbreak{}src/\allowbreak{}services/\allowbreak{}utils/\allowbreak{}FluxPeerManager.js}{32}. Liveness is a WebSocket ping
every \texttt{ws\allowbreak{}Ping\allowbreak{}Interval\allowbreak{}Ms} = 15{,}000\,ms; a peer that misses 3 consecutive pongs
(\texttt{ws\allowbreak{}Max\allowbreak{}Missed\allowbreak{}Pongs}) is force-terminated
\srccode{flux/\allowbreak{}ZelBack/\allowbreak{}config/\allowbreak{}default.js}{18--19}
\srccode{flux/\allowbreak{}ZelBack/\allowbreak{}src/\allowbreak{}services/\allowbreak{}utils/\allowbreak{}FluxPeerSocket.js}{86--145,209--212,298}. Placement itself
is gated on this liveness signal: the spawner runs only above an \texttt{app\allowbreak{}Sync\allowbreak{}Peer\allowbreak{}Threshold} of
12 live peers and pauses below an \texttt{app\allowbreak{}Sync\allowbreak{}Degraded\allowbreak{}Threshold} of 4
\srccode{flux/\allowbreak{}ZelBack/\allowbreak{}config/\allowbreak{}default.js}{268--269}
\srccode{flux/\allowbreak{}ZelBack/\allowbreak{}src/\allowbreak{}services/\allowbreak{}appLifecycle/\allowbreak{}appSpawner.js}{38--52}.

\subsection{Resource accounting}
\label{sec:fluxos-resources}

\begin{table}[H]
\centering
\caption{Per-tier resource allowances (total / available to applications).}
\begin{tabular}{l l l l l}
\toprule
Tier & CPU (units of 0.1 core) & RAM (MB) & HDD (GB) & Collateral, current / legacy (\flux) \\
\midrule
Cumulus ($\tC$) & 40 / 30 & 7{,}000 / 5{,}000 & 220 / 180 & 1{,}000 / 10{,}000 \\
Nimbus ($\tN$) & 80 / 70 & 30{,}000 / 28{,}000 & 440 / 400 & 12{,}500 / 25{,}000 \\
Stratus ($\tS$) & 160 / 150 & 61{,}000 / 59{,}000 & 880 / 840 & 40{,}000 / 100{,}000 \\
\bottomrule
\end{tabular}
\srccode{flux/\allowbreak{}ZelBack/\allowbreak{}config/\allowbreak{}default.js}{380--403}
\end{table}

Resources locked by the host system, subtracted before any application allocation:
\texttt{cpu}: 10 (one core, in units of 0.1 core), \texttt{ram}: 2{,}000\,MB, \texttt{hdd}: 60\,GB,
\texttt{extrahdd}: 20\,GB \srccode{flux/\allowbreak{}ZelBack/\allowbreak{}config/\allowbreak{}default.js}{375--378}, with a further 0.95
usable-disk multiplier \srccode{flux/\allowbreak{}ZelBack/\allowbreak{}src/\allowbreak{}services/\allowbreak{}appRequirements/\allowbreak{}hwRequirements.js}{385}.
\texttt{hwRequirements.\allowbreak{}checkAppHWRequirements} enforces, per component:
\begin{multline*}
\text{usable HDD} = 0.95\cdot\text{total HDD} - 60 - 20, \\
\text{reject if } \text{app HDD} > \text{usable HDD} - \text{HDD locked by other apps};
\end{multline*}
\begin{multline*}
\text{usable CPU (}\times 10\text{)} = 10\cdot\text{cores} - 10, \\
\text{reject if } 10\cdot\text{app CPU} > \text{usable CPU} - 10\cdot\text{CPU locked by other apps};
\end{multline*}
\begin{multline*}
\text{usable RAM} = \text{total RAM} - 2{,}000, \\
\text{reject if app RAM} > \text{usable RAM} - \text{RAM locked by other apps}
\end{multline*}
\srccode{flux/\allowbreak{}ZelBack/\allowbreak{}src/\allowbreak{}services/\allowbreak{}appRequirements/\allowbreak{}hwRequirements.js}{370--408}. A separate,
enterprise-only check additionally rejects an install that would leave 4 or fewer free vCores after
it (\texttt{free\allowbreak{}Cores\allowbreak{}After\allowbreak{}Install} = cores $-$ 1 system-reserved core $-$ locked cores $-$ app
cores), reserved so the CFS CPU-burst feature always has spare capacity to draw on
\srccode{flux/\allowbreak{}ZelBack/\allowbreak{}src/\allowbreak{}services/\allowbreak{}appRequirements/\allowbreak{}hwRequirements.js}{411--437}. That burst feature
is enabled by default (\texttt{cpuBurst.\allowbreak{}enabled: true}, CFS period \texttt{periodUs} = 100{,}000
(100\,ms), 1 reserved core) \srccode{flux/\allowbreak{}ZelBack/\allowbreak{}config/\allowbreak{}default.js}{440--453}, with the per-host
caveat, stated in the source itself, that "multiple burstable apps on the same host can
collectively oversubscribe during simultaneous spikes"
\srcdoc{flux/ZelBack/config/default.js:441-449}.

\subsection{Pricing}
\label{sec:fluxos-pricing}

Price is computed from a per-component resource vector $\resvec$, the requested instance count
$\ninst$, and the requested duration $\dur$ (in blocks, against $\blockslasting$), evaluated at the
transaction's block height $\hgt$ against a height-scheduled price table
(\texttt{config.\allowbreak{}fluxapps.\allowbreak{}price[]}, overridable by \texttt{chainmessages} of type \texttt{p} posted
on-chain by the Foundation multisig) \srccode{flux/\allowbreak{}ZelBack/\allowbreak{}src/\allowbreak{}services/\allowbreak{}utils/\allowbreak{}chainUtilities.js}{9--52}.
Concretely, for a \specv{4}-or-later (composed) app, \texttt{cpuPrice} $= \sum$(CPU units)
$\times\,\text{table.cpu}\times 10$; \texttt{ramPrice} $=\sum$(RAM units)
$\times\,\text{table.ram}/100$; \texttt{hddPrice} $=\sum$(HDD units) $\times\,\text{table.hdd}$;
plus \texttt{table.scope} if the app declares \texttt{nodes} or is \texttt{enterprise}, plus
\texttt{table.staticip} if it declares \texttt{staticip}, plus \texttt{table.port} per enterprise
port; that sum is quoted per three instances --- divided by three, ceiled to two decimals, and, from
height \num{1890000}, scaled by the requested instance count $\ninst$, with a floor of $0.50$ per
additional instance for the cheapest specifications
\srccode{flux/\allowbreak{}ZelBack/\allowbreak{}src/\allowbreak{}services/\allowbreak{}utils/\allowbreak{}appUtilities.js}{89--134}. The registration price is then
scaled by $\dur/\text{defaultExpire}$, where \texttt{defaultExpire} is $\blockslasting$ or
$4\times\blockslasting$ post-\PoN fork, re-ceiled, and only then floored at \texttt{minPrice}
\srccode{flux/\allowbreak{}ZelBack/\allowbreak{}src/\allowbreak{}services/\allowbreak{}appMessaging/\allowbreak{}messageVerifier.js}{733--743}; an update is priced
at 90\% of the fresh price minus a further discount proportional to the unused life of the prior
spec, floored again at \texttt{minPrice} \srccode{flux/\allowbreak{}ZelBack/\allowbreak{}src/\allowbreak{}services/\allowbreak{}appMessaging/\allowbreak{}messageVerifier.js}{801--816}.
\S16 gives the full pricing mechanics and treats the credits/renewal layer built on top of them;
here it is enough to record the one finding that changes what a payer can rely on:
\textbf{an underpaid registration or update is silently dropped, with only a server-side log
warning and no signal of any kind returned to the payer}
(\S\ref{sec:fluxos-lifecycle}; \srccode{flux/\allowbreak{}ZelBack/\allowbreak{}src/\allowbreak{}services/\allowbreak{}appMessaging/\allowbreak{}messageVerifier.js}{735--763,797--828}).
\S16 takes up what a correct failure signal would require.

\subsection{Moderation, as it exists today}
\label{sec:fluxos-moderation}

Docker image compliance is enforced against a GitHub-hosted list, not an on-chain or
network-quorum one: \texttt{get\allowbreak{}Blocked\allowbreak{}Repositores()} fetches
\texttt{helpers/\allowbreak{}blockedrepositories.\allowbreak{}json} from \texttt{raw.\allowbreak{}githubusercontent.\allowbreak{}com/\allowbreak{}RunOnFlux/\allowbreak{}flux},
cached in \texttt{fluxCaching.\allowbreak{}blockedRepositoriesCache}
\srccode{flux/\allowbreak{}ZelBack/\allowbreak{}src/\allowbreak{}services/\allowbreak{}appSecurity/\allowbreak{}imageManager.js}{179--195}; a parallel
\texttt{vettedrepositories.\allowbreak{}json} lets listed apps bypass user-defined blocked ports and repositories
\srccode{flux/\allowbreak{}ZelBack/\allowbreak{}src/\allowbreak{}services/\allowbreak{}appSecurity/\allowbreak{}imageManager.js}{202--236,395--410}.
\texttt{check\allowbreak{}Applications\allowbreak{}Compliance} runs every \texttt{image\allowbreak{}Compliance\allowbreak{}Interval\allowbreak{}Ms} =
3{,}600{,}000\,ms (1\,h), removing any locally installed app that matches the blocklist and waiting
3\,min between removals "so we don't have more removals at the same time"
\srccode{flux/\allowbreak{}ZelBack/\allowbreak{}src/\allowbreak{}services/\allowbreak{}serviceManager.js}{529--531}
\srccode{flux/\allowbreak{}ZelBack/\allowbreak{}config/\allowbreak{}default.js}{319}
\srccode{flux/\allowbreak{}ZelBack/\allowbreak{}src/\allowbreak{}services/\allowbreak{}appSecurity/\allowbreak{}imageManager.js}{644--676}. The equivalent
GitHub-sourced enterprise node/owner map resyncs every 6\,h with fetch failure explicitly swallowed
by a \texttt{.catch} handler, proceeding on the last-cached value
\srccode{flux/\allowbreak{}ZelBack/\allowbreak{}src/\allowbreak{}services/\allowbreak{}serviceManager.js}{130--136}; the evidence gathered for this
paper did not separately trace a distinct failure-handling code path for the blocklist fetch itself
beyond its cache, but the two fetches share the same GitHub-hosted, silently-cached design.
\S23 analyzes what this means for censorship-resistance and for the roadmap's stated network-quorum
vision; this section states the mechanism and defers the analysis.

\subsection{Trust and failure model}
\label{sec:fluxos-trust}

\begin{enumerate}
\item \textbf{Placement convergence assumes peer honesty.} Instance-count convergence depends
  entirely on nodes honestly reporting and honoring \texttt{broadcastedAt} rank; nothing in the
  paths read re-verifies it against the signed envelope's own timestamp for this message type
  \srcinferred (\S\ref{sec:fluxos-placement}).
\item \textbf{Underpayment is swallowed, not rejected.} A registration or update that pays too
  little produces only a log warning; the payer receives no on-chain refund and no distinguishable
  failure signal \srccode{flux/\allowbreak{}ZelBack/\allowbreak{}src/\allowbreak{}services/\allowbreak{}appMessaging/\allowbreak{}messageVerifier.js}{735--763,797--828}.
\item \textbf{GitHub is the moderation trust root.} The image blocklist, vetted-repository list,
  and enterprise owner map are all fetched from one GitHub repository on a schedule, with fetch
  failure silently swallowed and the last-good value kept
  \srccode{flux/\allowbreak{}ZelBack/\allowbreak{}src/\allowbreak{}services/\allowbreak{}serviceManager.js}{130--136}; a GitHub outage or a compromise
  of that specific file is a single point of moderation trust for the whole network, not a
  network-quorum mechanism (\S\ref{sec:fluxos-moderation}).
\item \textbf{Boot is a flat-backoff infinite retry, not a circuit breaker.} Any uncaught exception
  in the boot function restarts it after a fixed 15{,}000\,ms with no retry cap and no escalation,
  so a persistently failing dependency retries forever rather than surfacing a fatal state
  \srccode{flux/\allowbreak{}ZelBack/\allowbreak{}src/\allowbreak{}services/\allowbreak{}serviceManager.js}{568--572}.
\item \textbf{Distributed enforcement trusts a synchronized local view.} The node-status monitor's
  network-wide eviction pass — which every confirmed node runs independently against every other
  node's claimed locations — trusts its own local copy of the deterministic node list; if that
  local view lags, it can evict still-valid locations or fail to evict dead ones \srcinferred.
\item \textbf{Self-correction can amplify churn under partition.} Because every node independently
  re-derives whether an application is over-replicated and removes itself on that basis, a
  partition that produces disjoint \texttt{runningSince} orderings can cause more nodes to remove
  an application than intended, or none, until gossip re-converges; there is no distributed lock
  (\S\ref{sec:fluxos-placement}, remark above).
\item \textbf{\texttt{admin} is a local file, not a bound credential.} \texttt{admin} privilege is
  granted to whoever controls \texttt{userconfig.\allowbreak{}initial.\allowbreak{}zelid} on that specific node's disk; there
  is no additional hardware or HSM binding checked at this privilege level
  \srccode{flux/\allowbreak{}ZelBack/\allowbreak{}src/\allowbreak{}services/\allowbreak{}verificationHelperUtils.js}{19--46}.
\end{enumerate}

%% Drain this section's float queue before the next begins.
\FloatBarrier
